\section*{अध्याय \ref{ch:b1:chemical-biogenic-sediments} --- कार्बोनेट, सिलिसियस और वाष्पनिक तलछट}

\begin{solution}{exo:b1:chemical-biogenic-sediments:1}
क्लोराइड $0.5503\times35.165 = \qty{19.35}{g}$; सोडियम $0.3066\times35.165 =
\qty{10.78}{g}$; सल्फ़ेट \qty{2.711}{g}; मैग्नीशियम \qty{1.284}{g}। सोडियम और क्लोराइड मिलकर $0.5503 + 0.3066 = 0.857$ हैं, यानी लवण का \qty{85.7}{\%}।
\end{solution}

\begin{solution}{exo:b1:chemical-biogenic-sediments:2}
बचा हुआ अंश $1/\mathrm{CF}$ है: कैल्शियम कार्बोनेट $1/1.8 = 0.556$ बचा, \qty{44.4}{\%} वाष्पित; जिप्सम $0.263$ बचा, \qty{73.7}{\%} वाष्पित; हैलाइट $0.094$ बचा, \qty{90.6}{\%} वाष्पित; मैग्नीशियम सल्फ़ेट $0.014$ बचा, \qty{98.6}{\%} वाष्पित।
\end{solution}

\begin{solution}{exo:b1:chemical-biogenic-sediments:3}
(क)~$\rho_s(1-\phi) = G/a = 8.4/0.0042 = \qty{2000}{kg.m^{-3}}$। (ख)~ऐरागोनाइट के साथ, $1 - 2000/2930 = 0.32$, यानी \qty{32}{\%} छिद्रता; कैल्साइट के साथ, $1 - 2000/2710 = 0.26$। (ग)~$a = 3.7/2000 = \qty{1.85e-3}{m}$, यानी प्रति वर्ष \qty{1.85}{mm}।
\end{solution}

\begin{solution}{exo:b1:chemical-biogenic-sediments:4}
\qty{4000}{m} तक मात्रा केवल धीरे गिरती है (88 से 80); \qty{4000}{m} के नीचे वह तीव्रता से गिरती है: \omterm{def:b1:chemical-biogenic-sediments:ccd}{लाइसोक्लाइन} लगभग \qty{4000}{m} पर है। मात्रा \qty{4700}{m} (12) और \qty{4800}{m} (6) के बीच \qty{10}{\%} से गुज़रती है: $4700 + 100\times(12-10)/(12-6) = \qty{4733}{m}$, यानी CCD लगभग \qty{4730}{m} पर, \omterm{def:b1:chemical-biogenic-sediments:ccd}{लाइसोक्लाइन} से लगभग \qty{700}{m} नीचे।
\end{solution}

\begin{solution}{exo:b1:chemical-biogenic-sediments:5}
(क)~\ce{Ca^2+ + 2 HCO3- -> CaCO3 + CO2 + H2O}: एक मोल \ce{CO2} प्रति मोल \ce{CaCO3}, इसलिए $44.01/100.09 = \qty{0.440}{kg}$ \ce{CO2} प्रति किलोग्राम \ce{CaCO3}। (ख)~$3.7\times0.440 = \qty{1.63}{kg}$ \ce{CO2} प्रति वर्ग मीटर और प्रति वर्ष। (ग)~प्रयुक्त दो हाइड्रोजन कार्बोनेट आयनों में से केवल एक कार्बन \omterm{def:b1:rock-forming-minerals:mineral}{खनिज} में जाता है; दूसरा कार्बन डाइऑक्साइड के रूप में निकलता है। कैल्सीकरण पानी से क्षारकता निकालता है और उसकी घुले कार्बन को धारण करने की क्षमता घटाता है; कितना \ce{CO2} वायु तक पहुँचता है, यह समुद्री जल की बफ़रिंग पर निर्भर है, जो रसायन-विज्ञान की पुस्तकों के कार्बोनेट साम्यों का विषय है।
\end{solution}

\begin{solution}{exo:b1:chemical-biogenic-sediments:6}
(क)~$t_c = (2000/350)^2 = \qty{32.7}{Myr}$; \qty{25}{km} प्रति दस लाख वर्ष पर, $25\times32.7 = \qty{816}{km}$ अक्ष से दूर। (ख)~$t_c = (2600/350)^2 =
\qty{55.2}{Myr}$, \qty{1380}{km}। (ग)~\qty{60}{Myr} पुराने तल ने पहले मामले में $60 - 32.7 = \qty{27.3}{Myr}$ तक लाल मृत्तिका पाई है, और दूसरे में केवल $60 - 55.2 = \qty{4.8}{Myr}$ तक।
\end{solution}

\begin{solution}{exo:b1:chemical-biogenic-sediments:7}
(क)~हाँ: $a = 4.2 > r = 3$, और \omterm{def:b1:chemical-biogenic-sediments:reef}{भित्ति} समुद्र पर \qty{1.2}{mm} प्रति वर्ष लाभ करती है जब तक उसका शीर्ष सतह तक न पहुँचे। (ख)~$D$ हर वर्ष $10 - 4.2 = \qty{5.8}{mm}$ बढ़ता है; \qty{2}{m} से \qty{50}{m} तक जाने में $48/0.0058 = \num{8300}$ वर्ष लगते हैं। (ग)~क्योंकि जैसे-जैसे पानी गहरा होता है, उसकी वृद्धि दर घटती है: प्रवालों में रहने वाले और उन्हें पोषण देने वाले शैवालों तक कम प्रकाश पहुँचता है, इसलिए $a$ घटता है और \omterm{def:b1:chemical-biogenic-sediments:reef}{भित्ति} और तेज़ी से पिछड़ती जाती है।
\end{solution}

\begin{solution}{exo:b1:chemical-biogenic-sediments:8}
(क)~अभिक्रिया है \ce{CaSO4.2H2O -> CaSO4 + 2 H2O}। (ख)~मोलर आयतन हैं $172.17/2.30 =
\qty{74.9}{cm^3}$ और $136.14/2.97 = \qty{45.8}{cm^3}$, इसलिए ठोस अपने आयतन का $1 - 45.8/74.9 = \qty{38.8}{\%}$ खोता है। (ग)~$2\times18.02/172.17 = 0.209$: \qty{20.9}{\%} द्रव्यमान जल के रूप में निकलता है। गहराई पर निकला यह जल आसपास की चट्टानों में छिद्र-दाब बढ़ाता है; एक सील के नीचे वह चट्टानों को कमज़ोर कर सकता है और भ्रंशों को सरकने में मदद कर सकता है, और वह अन्य लवणों को घोल और ढो सकता है।
\end{solution}

\begin{solution}{exo:b1:chemical-biogenic-sediments:9}
(क)~$9.2/255 = \qty{3.6}{\%}$ संरक्षित, \qty{96.4}{\%} घुला। (ख)~जहाँ उत्पादन बहुत अधिक है, उसका कुछ प्रतिशत भी एक बड़ा प्रवाह है, और जहाँ थोड़ी मृत्तिका या कार्बोनेट उसे तनु करता है: दक्षिणी महासागर, उत्तरी प्रशांत और भूमध्यरेखीय प्रशांत के ऊपर उठते जल के नीचे, और CCD के नीचे। (ग)~पतले, नाज़ुक फ़स्टूल पहले घुलते हैं और मज़बूत बचते हैं: तलछट सबसे कठोर प्रजातियों का चयन रखता है, उस समुदाय को नहीं जो पानी में रहता था।
\end{solution}

\begin{solution}{exo:b1:chemical-biogenic-sediments:10}
(क)~$EA = 1.5\times\qty{1e10}{m^2} = \qty{1.5e10}{m^3}$ प्रति वर्ष। (ख)~$Q_{\mathrm{in}} = \qty{1.8e10}{m^3}$, $Q_{\mathrm{out}} = \qty{0.3e10}{m^3}$ और $k = 6$ के साथ: जिप्सम की देहली (3.8) से आगे, पर हैलाइट की (10.6) से नीचे, इसलिए कैल्शियम कार्बोनेट और जिप्सम अवक्षेपित होते हैं। $Q_{\mathrm{in}} =
\qty{1.575e10}{m^3}$, $Q_{\mathrm{out}} = \qty{0.075e10}{m^3}$ और $k = 21$ के साथ, मैग्नीशियम जैसे एक आयन के लिए जो घोल में रहता है: ब्राइन हैलाइट की देहली पार करती है, इसलिए हैलाइट भी अवक्षेपित होता है, पर मैग्नीशियम लवण (70) नहीं। (ग)~हैलाइट के लिए $Q_{\mathrm{in}}/(Q_{\mathrm{in}} - EA) \ge 10.6$ आवश्यक है, अर्थात् $Q_{\mathrm{in}} \le 10.6\,EA/9.6 = 1.104\,EA = \qty{1.66e10}{m^3}$ प्रति वर्ष: अंतःप्रवाह वाष्पन से अधिक से अधिक लगभग एक-दसवाँ भाग हो सकता है।
\end{solution}

\begin{solution}{exo:b1:chemical-biogenic-sediments:11}
(क)~समुद्री जल में प्रति किलोग्राम \qty{67.3}{mg} ब्रोमाइड है; 10.6~गुना पर, ब्राइन में $67.3\times10.6 = \qty{713}{mg.kg^{-1}}$ और पहले हैलाइट में $0.032\times713 = \qty{22.8}{mg.kg^{-1}}$। (ख)~30~गुना पर, $67.3\times30 =
\qty{2019}{mg.kg^{-1}}$ ब्राइन में और हैलाइट में \qty{64.6}{mg.kg^{-1}}। (ग)~हैलाइट ब्राइन की सांद्रता का केवल 3.2~सौवें भाग लेता है, इसलिए लगभग सारा ब्रोमाइड घोल में रहता है और पानी के वाष्पित होने के साथ सांद्र होता जाता है। (घ)~वृद्धि ब्राइन की सांद्रता अंकित करती है; अचानक गिरावट उसके तनुकरण को अंकित करती है, द्रोणी में समुद्री जल के नए अंतःप्रवाह या ताज़े पानी द्वारा, इससे पहले कि हैलाइट बढ़ता रहे।
\end{solution}

\begin{solution}{exo:b1:chemical-biogenic-sediments:12}
(क)~$t_c = \qty{32.65}{Myr}$, इसलिए $0.02\times\num{32650} = \qty{653}{m}$ ऊज़ (\num{32650} हज़ार-वर्षीय अंतराल, प्रत्येक \qty{2}{cm} पर) और $0.001\times\num{37350} = \qty{37}{m}$ लाल मृत्तिका। (ख)~$t_c = (2500/350)^2 =
\qty{51.0}{Myr}$: \qty{1020}{m} ऊज़ और \qty{19}{m} लाल मृत्तिका। गहरा CCD कहीं अधिक मोटा ऊज़ देता है, और गहरे अटलांटिक तल की लाल मृत्तिका उसी आयु के प्रशांत की तुलना में पतली है। (ग)~CCD कुछ समय के लिए और गहरा रहा होगा, इसलिए स्थान उसके ऊपर फिर उठ गया; या प्लेट ने स्थान को एक उच्च उत्पादकता की पेटी के नीचे पहुँचाया, जैसे भूमध्य रेखा, जहाँ अधिक कार्बोनेट बरसा; या परत किसी ढलान-खिसकन या टर्बिडिटी धारा से समुद्र-तल के ऊँचे भाग से लाई गई।
\end{solution}

\begin{solution}{pb:b1:chemical-biogenic-sediments:1}
\textbf{1.} $1030\times1000 = \qty{1.03e6}{kg}$ समुद्री जल का, जिसमें
$1.03\times10^6\times0.035165 = \qty{36200}{kg}$ लवण है।

\textbf{2.} जिप्सम $1.03\times10^6\times1.581\times10^{-3} = \qty{1628}{kg}$;
हैलाइट $1.03\times10^6\times0.02741 = \qty{28230}{kg}$।

\textbf{3.} कैल्शियम कार्बोनेट $113.3/2710 = \qty{0.042}{m}$; जिप्सम $1628/2300 =
\qty{0.71}{m}$; हैलाइट $28230/2170 = \qty{13.0}{m}$।

\textbf{4.} एनहाइड्राइट का द्रव्यमान $1628\times136.14/172.17 =
\qty{1288}{kg}$ है और मोटाई $1288/2970 = \qty{0.43}{m}$।

\textbf{5.} सूखा नमक $0.110 + 1.250 + 27.41 + 6.34 = \qty{35.11}{g}$ प्रति किलोग्राम है, जिसमें हैलाइट $27.41/35.11 = \qty{78}{\%}$ है: सोडियम और क्लोराइड अधिकांश घुले लवण बनाते हैं, और हैलाइट लगभग सब कुछ ले लेता है।

\textbf{6.} नीचे से ऊपर: कैल्शियम कार्बोनेट, जिप्सम, हैलाइट, बिटर्न लवण। हर \omterm{def:b1:rock-forming-minerals:mineral}{खनिज} तब अवक्षेपित होना शुरू करता है जब ब्राइन उसकी संतृप्ति तक पहुँचती है, सबसे कम विलेय पहले, और पहले जमे हुए पर बैठता है।

\textbf{7.} ब्राइन की गहराई $1000/\mathrm{CF}$ है: कैल्शियम कार्बोनेट के लिए \qty{556}{m}, जिप्सम के लिए \qty{263}{m}, हैलाइट के लिए \qty{94}{m}।

\textbf{8.} $1 - 1/10.6 = \qty{90.6}{\%}$।

\textbf{9.} मैग्नीशियम सल्फ़ेट के लिए $1000/70 = \qty{14}{m}$, पोटैशियम लवणों के लिए $1000/90 = \qty{11}{m}$।

\textbf{10.} कैल्शियम कार्बोनेट बनने पर लगभग \qty{1.10}{g.cm^{-3}}, जिप्सम के साथ 1.13, हैलाइट के साथ 1.22 और बिटर्न लवणों के साथ 1.30 से अधिक, समुद्री जल के 1.03 के मुकाबले। भीतर बहता समुद्री जल सघन ब्राइन पर तैरता है, जो तल पर रहती है; अवरुद्ध द्रोणी में ब्राइन अंतःप्रवाह के नीचे महासागर को लौट सकती है।

\textbf{11.} वह केवल \qty{4}{cm} मोटी है, \qty{13}{m} हैलाइट के मुकाबले; और वह किनारों पर बनती है, जहाँ आने वाला पानी पहले सांद्र होता है, केंद्र से अधिक, इसलिए वह जिप्सम के आधार में मिल जाती है या उसकी जगह ले लेती है।

\textbf{12.} हैलाइट ब्राइन के ब्रोमाइड का केवल छोटा भाग लेता है (अभ्यास~\ref{exo:b1:chemical-biogenic-sediments:11}), इसलिए हैलाइट बढ़ने के साथ ब्राइन ब्रोमाइड में और समृद्ध होती है, और हैलाइट की हर नई परत और अधिक लेती है।

\textbf{13.} $1500/13.01 = 115.3$: 116 भराव।

\textbf{14.} $116\times0.434 = \qty{50}{m}$ एनहाइड्राइट, एक परत के रूप में नहीं बल्कि 116 परतों के रूप में, हर एक आधा मीटर से कुछ कम मोटी, अपने भराव की हैलाइट के आधार पर: एनहाइड्राइट और हैलाइट का एक चक्रीय अनुक्रम।

\textbf{15.} प्रति भराव, $1.03\times10^6\times6.34\times10^{-3}/2000 =
\qty{3.3}{m}$, इसलिए 116 भरावों के लिए \qty{379}{m}। वे सबसे विलेय लवण हैं: समुद्री जल का अगला अंतःप्रवाह उन्हें फिर से घोल देता है, और वह द्रोणी जिसकी ब्राइन अंत से पहले बह जाती है, उन तक कभी नहीं पहुँचती।

\textbf{16.} हर पुनर्भराव पर ब्राइन तनु होती है: नई हैलाइट के आधार पर ब्रोमाइड की मात्रा तेज़ी से गिरती है और ऊपर की ओर फिर बढ़ती है, एक आरी-दाँत जैसी रूपरेखा, हर भराव के लिए एक दाँत।

\textbf{17.} $2000/13.01 = 154$ भराव।

\textbf{18.} $5000\times365.25\times86400 = \qty{1.578e11}{m^3}$ प्रति वर्ष, जो $\qty{1.0e11}{m^2}$ पर फैला है: $E = \qty{1.58}{m}$ प्रति वर्ष।

\textbf{19.} $1.578\times10^{11}\times1030\times0.035165 = \qty{5.72e12}{kg}$ लवण प्रति वर्ष।

\textbf{20.} $1.578\times10^{11}\times1030\times0.02741/2170 =
\qty{2.05e9}{m^3}$ हैलाइट प्रति वर्ष, यानी \qty{0.0205}{m}, तल पर लगभग \qty{2}{cm} प्रति वर्ष।

\textbf{21.} $1500/0.0205 = \num{73000}$ वर्ष।

\textbf{22.} स्तर \qty{1.58}{m} प्रति वर्ष गिरता है; \qty{1000}{m} पानी $1000/1.578 = 634$ वर्षों में चला जाता है।

\textbf{23.} हैलाइट तब प्रकट होता है जब पानी का \qty{90.6}{\%} जा चुका हो, $905.7/1.578 = 574$ वर्ष बाद।

\textbf{24.} \qty{1.5}{km} हैलाइट में द्रोणी के 116 भरावों का नमक लगा (115.3, ऊपर की ओर पूर्णांकित), चाहे वे एक के बाद एक आए हों जब द्रोणी सूखती और फिर भरती थी, या देहली के ऊपर से लगातार, जैसा प्रश्न~21 में है, समुद्री जल की एक द्रोणी हर $1000/1.578 = 634$ वर्ष में।
\end{solution}
